The within-country tests in Sections~\ref{sec:test_h} and~\ref{sec:test_i} ask whether Belgian buyers reallocate purchases away from regulated \emph{domestic} suppliers. The natural complementary question is whether buyers reallocate away from regulated \emph{foreign} suppliers --- the canonical cross-border leakage channel that motivates CBAM and a large fraction of the trade-policy environmental literature. We test this on the Belgian customs panel using the identification strategy of \citet{CosterMejeanDiGiovanni2024} (CMdG hereafter), which they implement on French firm-level customs data and find a $\sim$12 percentage-point Phase~3 reallocation of regulated-product imports toward non-ETS source countries. Replicating their specification on Belgium yields a result that is two orders of magnitude smaller than the French estimate and points in the opposite direction.

\subsection{Specification}\label{sec:across:spec}

Following \citet{CosterMejeanDiGiovanni2024} Equation~(1), we estimate
\begin{equation}
    y_{f,p,i,t} \;=\; \sum_{\phi \in \{1, 2, 3\}} \beta_\phi \,\cdot\, \mathbf{1}[\text{regulated}]_p \,\times\, \mathbf{1}[t \in \text{Phase } \phi] \;+\; \alpha_{f,p,i} \;+\; \delta_{i,t} \;+\; \delta_{s,t} \;+\; \varepsilon_{f,p,i,t}
    \label{eq:cmdg_spec}
\end{equation}
on a balanced panel of (Belgian buyer firm $f$, HS6 product $p$, source country $i$, year $t$) cells over 2000--2019, restricted to imports from non-ETS source countries (per CMdG p.~12--13, who restrict to avoid contamination from intra-EU Intrastat threshold changes). $\mathbf{1}[\text{regulated}]_p$ flags HS6 codes that map (via the Phase~0 concordance bridge described in the appendix) to NACE 4-digit categories containing at least one ETS installation. $\alpha_{f,p,i}$ is a firm $\times$ product $\times$ country fixed effect, $\delta_{i,t}$ a country $\times$ year fixed effect, and $\delta_{s,t}$ a (buyer-)sector $\times$ year fixed effect. Phase~1 (2005--2008), Phase~2 (2009--2012), and Phase~3 (2013--2019) span the EU ETS post-period; the implicit reference window is 2000--2004.

We estimate equation~\eqref{eq:cmdg_spec} with two outcomes:
\begin{align*}
    \text{Panel A — share:} \quad y_{f,p,i,t} &\;=\; \frac{\text{value}_{f,p,i,t}}{\sum_{p',i' \in \text{non-ETS}} \text{value}_{f,p',i',t}}, \\
    \text{Panel B — probability:} \quad y_{f,p,i,t} &\;=\; \mathbf{1}[\text{value}_{f,p,i,t} > 0].
\end{align*}
Standard errors are two-way clustered on firm and on country. We report the same six fixed-effect column structure as \citet{CosterMejeanDiGiovanni2024} Table~1, with column~(5) (firm $\times$ product $\times$ country, country $\times$ year, sector $\times$ year) as the preferred specification.

\subsection{Sample}\label{sec:across:sample}

Starting from the raw NBB customs records, we apply the following filter cascade to match CMdG's identifying sample as closely as possible: imports only (flow $=$ I); manufacturing buyer (NACE 2-digit $\in$ \{10--33\}); buyer in a regulated-intensive sector (i.e., a NACE 2-digit code where the IO matrix shows non-trivial use of regulated upstream inputs); core-input filter (the imported product is in the buyer's top decile of upstream input intensity in the Belgian Use table); and dropping capital goods via the BEC concordance. After balancing the (firm $\times$ HS6 $\times$ country $\times$ year) panel with zero-fill and restricting to non-ETS source countries, we obtain 8.6~million observations on 3{,}000--5{,}000 distinct buyer firms, 19 buyer NACE 2-digit sectors, and 199 partner countries; 35.4\% of cells are regulated. The CMdG France sample is 7.5~million observations on $\sim$27{,}000 firms; Belgium has fewer firms but comparable panel breadth per firm.

\subsection{Headline}\label{sec:across:headline}

Table~\ref{tab:cmdg_belgium} reports the CMdG-exact column~(5) coefficients for Belgium and the corresponding French coefficients from \citet{CosterMejeanDiGiovanni2024} Table~1.

\begin{table}[H]
    \centering
    \caption{CMdG specification on the Belgian customs panel, column~(5) preferred fixed effects (firm $\times$ product $\times$ country, country $\times$ year, sector $\times$ year). Two-way clustered standard errors. France comparison from \citet{CosterMejeanDiGiovanni2024} Table~1.}
    \label{tab:cmdg_belgium}
    \footnotesize
    \begin{tabular}{lcccc}
        \toprule
        & \multicolumn{2}{c}{\textbf{Belgium (this paper)}} & \multicolumn{2}{c}{\textbf{France (CMdG Table 1)}} \\
        \cmidrule(lr){2-3} \cmidrule(lr){4-5}
        Coefficient & $\widehat{\beta}$ (s.e.) & $p$-value & $\widehat{\beta}$ & sig. \\
        \midrule
        \multicolumn{5}{l}{\emph{Panel A — share}} \\
        Phase 1 (2005--08) $\times \mathbf{1}[\text{reg}]$    & $-0.0014$ (0.0008) & 0.091    & ---       & ---     \\
        Phase 2 (2009--12) $\times \mathbf{1}[\text{reg}]$    & $-0.0037$ (0.0011) & 0.001*** & $+0.110$  & ***     \\
        Phase 3 (2013--19) $\times \mathbf{1}[\text{reg}]$    & $-0.0024$ (0.0009) & 0.007*** & $+0.121$  & ***     \\
        \midrule
        \multicolumn{5}{l}{\emph{Panel B — probability}} \\
        Phase 1 $\times \mathbf{1}[\text{reg}]$               & $-0.0041$ (0.0051) & 0.43     & ---       & ---     \\
        Phase 2 $\times \mathbf{1}[\text{reg}]$               & $-0.0128$ (0.0095) & 0.18     & (positive)& ---     \\
        Phase 3 $\times \mathbf{1}[\text{reg}]$               & $+0.0041$ (0.0097) & 0.67     & $+0.071$  & ***     \\
        \bottomrule
    \end{tabular}
\end{table}

The three immediate readings of Table~\ref{tab:cmdg_belgium}:

\begin{itemize}
    \item \textbf{Belgium does not exhibit the cross-border leakage CMdG document for France.} The Phase~3 share coefficient is $-0.0024$ for Belgium versus $+0.121$ for France --- two orders of magnitude smaller and pointing in the opposite direction.

    \item \textbf{Where statistically significant in Belgium, the share coefficients are wrong-signed for leakage.} Phases~2 and~3 both deliver \emph{negative} share coefficients at conventional confidence levels ($p = 0.001$ and $p = 0.007$ respectively). Belgian buyers' share of imports going to (regulated $\times$ non-ETS) cells \emph{fell} relative to (unregulated $\times$ non-ETS) cells under the post-2005 ETS regime --- the opposite of what the leakage hypothesis predicts.

    \item \textbf{The probability margin is null in Belgium.} The extensive-margin coefficient on Phase~3 is $+0.0041$, two orders of magnitude smaller than the French $+0.071$, and not statistically distinguishable from zero. We discuss the non-zero pre-trend on the probability margin in Section~\ref{sec:across:eventstudy}, which complicates a literal interpretation of the level coefficients on Panel B.
\end{itemize}

The cross-column pattern across all six fixed-effect specifications --- not reported in Table~\ref{tab:cmdg_belgium} but available in the appendix --- is consistent with the column~(5) headline: share coefficients fall in the $[-0.005, +0.002]$ range, probability coefficients in $[-0.02, +0.01]$. No specification recovers CMdG-style positive coefficients of magnitude approaching $+0.1$.

\subsection{Robustness across control-group choices}\label{sec:across:robustness}

To check that the headline does not depend on a specific choice of treatment-vs-control comparison, we re-run equation~\eqref{eq:cmdg_spec} under four alternative control-group specifications, holding the column~(5) fixed effects constant. Table~\ref{tab:cmdg_robust} summarizes.

\begin{table}[H]
    \centering
    \caption{Robustness across control-group choices, Phase~3 coefficients on the Belgian customs panel. Columns~(A)--(D) vary the sample and the implicit comparison cell; column~(5) FE structure throughout.}
    \label{tab:cmdg_robust}
    \footnotesize
    \begin{tabular}{lllcc}
        \toprule
        Spec & Sample & Treatment vs.\ control & Panel A (share) & Panel B (prob.) \\
        \midrule
        \textbf{A} (CMdG baseline) & non-ETS countries only & regulated vs.\ unregulated & $-0.0024$** & $+0.0041$ \\
        \textbf{B} (CMdG Fig.\ 4 conceptual) & regulated products only & non-ETS vs.\ ETS countries & broken & broken \\
        \textbf{C} & full sample & (reg $\times$ non-ETS) vs.\ pooled rest & $-0.0010$ & $-0.0054$ \\
        \textbf{D} & full sample & triple-difference (no CMdG analog) & $-0.0010$ & $-0.0153^{\dagger}$ \\
        \bottomrule
        \multicolumn{5}{l}{\footnotesize $^{\dagger}$ marginal at $p \approx 0.10$. ``Broken'' = VCOV not positive definite.}
    \end{tabular}
\end{table}

The take-aways:

\begin{enumerate}
    \item \textbf{Spec~B is uninterpretable in the Belgian data.} The (regulated $\times$ ETS-country) cell that Spec~B uses as its implicit control is too thin in the Belgian customs panel: once we restrict to regulated products only and identify off non-ETS-vs-ETS country variation, the variance--covariance matrix fails to be positive definite. CMdG's Figure~4 conceptual analog cannot be cleanly estimated on Belgium.
    \item \textbf{Specs~A, C, D agree directionally.} All three converging specifications return small null-to-negative coefficients in Phase~3 (and stronger negative coefficients in Phase~2; not shown but reported in the IMPORT\_LEAKAGE companion document). The headline does not depend on which control we choose.
    \item \textbf{Phase~2 carries the strongest negative signal.} In specs C and D, the Phase~2 share coefficient is $-0.0014$** and the probability coefficient $-0.0291$*** --- economically modest but statistically distinct from zero. These specifications coincide with the post-2008 financial crisis and EU sovereign debt crisis aftermath, and are the periods in which Belgium's regulated $\times$ non-ETS share fell most sharply relative to the alternative control-group choices.
    \item \textbf{No specification recovers a positive headline.} The leakage CMdG identify for France is absent in the Belgian customs panel under any reasonable identification.
\end{enumerate}

\subsection{Event study}\label{sec:across:eventstudy}

Figure~\ref{fig:cmdg_event_study} plots the year-by-year coefficients from a continuous-time analog of equation~\eqref{eq:cmdg_spec}, with reference year $t = 2004$ (the last pre-EU-ETS year).

\begin{figure}[H]
    \centering
    \includegraphics[width=0.85\textwidth]{../output/figures/phase2_cmdj_figure3.png}
    \caption{Year-by-year coefficients for the Belgian CMdG specification, full sample, column~(5) fixed effects. Reference year 2004. Top panel: share outcome. Bottom panel: probability outcome.}
    \label{fig:cmdg_event_study}
\end{figure}

The pre-trend (2000--2003) is informative on the two panels:
\begin{itemize}
    \item \textbf{Panel A (share): clean pre-trend.} Coefficients in 2000--2003 are essentially zero ($-0.0002$ to $-0.0001$, all 95\% confidence intervals containing zero). Identification under parallel-trends is well-supported on the share outcome.
    \item \textbf{Panel B (probability): non-zero pre-trend.} Coefficients in 2000--2003 are negative at $-0.012$ to $-0.014$ with 95\% confidence intervals excluding zero. Treated cells were already structurally less likely to be active than control cells \emph{before} any ETS treatment.
\end{itemize}

The Panel~B pre-trend complicates a clean interpretation of the level coefficients on the probability outcome: we cannot rule out that the post-2005 widening to $-0.026$ to $-0.032$ (in the 2008--2014 window, before partial recovery toward zero by 2017--2019) is a continuation of the pre-existing trend rather than an ETS-induced effect. The Panel~A pre-trend is clean and the share-outcome conclusions are not subject to this concern. CMdG's France data shows a flat pre-trend on both panels, so this complication is specific to Belgium.

The pre-trend has a plausible structural explanation: regulated HS6 products are concentrated in chemicals, basic metals, and refining (HS chapters 27, 28, 29, 72--76), which are commodity-like and dominated by a handful of specialised sourcing relationships per buyer firm. Regulated cells therefore have a lower density of (firm $\times$ product $\times$ country) triplets to begin with, mechanically lowering the extensive-margin coefficient even absent any ETS effect.

\subsection{Belgium versus France}\label{sec:across:bevfrance}

Belgium and France implement the same EU ETS regime under the same regulatory architecture and the same Phase~III/IV price path. They differ in industrial composition and trade integration. Table~\ref{tab:bevfrance} summarises the comparison.

\begin{table}[H]
    \centering
    \caption{Belgium vs.\ France on the CMdG specification.}
    \label{tab:bevfrance}
    \footnotesize
    \begin{tabular}{lcc}
        \toprule
        & Belgium (this paper) & France \citep{CosterMejeanDiGiovanni2024} \\
        \midrule
        Sample period                  & 2000--2019                & 2000--2019 \\
        Distinct firms                 & 3{,}000--5{,}000          & 27{,}000 ($\sim$6$\times$ larger) \\
        Total panel rows               & 8.6M                      & 7.5M \\
        Buyer NACE 2-digit sectors     & 19                        & 20 (NAF-138 $\to$ NACE 2d) \\
        Phase~2 share, col.\ (5)       & $-0.0037$***              & $+0.110$*** \\
        Phase~3 share, col.\ (5)       & $-0.0024$**               & $+0.121$*** \\
        Phase~3 prob., col.\ (5)       & $+0.004$ (n.s.)           & $+0.071$*** \\
        Pre-trend (probability)        & $-0.012$ to $-0.014$ (non-zero) & flat \\
        Sign of headline               & wrong (negative)          & right (positive) \\
        \bottomrule
    \end{tabular}
\end{table}

Three working hypotheses for the divergence, in our preferred order of plausibility:
\begin{enumerate}
    \item \emph{Intra-EU substitution as the reaction margin.} Belgian buyers post-2005 may have shifted regulated-product sourcing toward EU-domestic suppliers (also subject to the EU ETS, but cheaper to reach logistically and free of currency or customs friction) rather than toward non-ETS countries. The negative Phase~2 and Phase~3 coefficients on Belgium are consistent with this: the regulated $\times$ non-ETS share fell because Belgian buyers consolidated to EU sources for the regulated inputs. France, with lower intra-EU trade integration, has more headroom to substitute toward non-EU sources at the margin.

    \item \emph{Pre-existing chemicals saturation.} The Antwerp chemicals cluster competes with Chinese, Indian, and Saudi chemical exporters on global markets and was already heavily integrated with non-ETS sources before 2005. Belgian buyers of regulated chemical inputs may have sourced as much as they were ever going to from non-ETS countries pre-treatment; further shifts post-treatment hit a saturation ceiling. France's chemical sector is smaller in relative terms, leaving more room for post-treatment reallocation.

    \item \emph{Methodological --- sample-size attenuation in column~(5).} With $\sim$3{,}000--5{,}000 Belgian firms versus France's 27{,}000, the firm $\times$ product $\times$ country fixed-effect block is sparser in Belgium, and the country $\times$ year fixed effects absorb a larger share of identifying variation. Cross-column comparison shows that the columns without country $\times$ year fixed effects (cols.\ 1, 2, 4, 6) deliver share coefficients closer to $\pm 0.005$ --- still small in absolute terms, but at least directionally consistent with France in some cells. Column~(5) absorbs whatever Belgian-leakage signal exists into the FE structure.
\end{enumerate}

We cannot separately identify these three channels with the current data. Hypotheses~(1) and~(2) are substantive (Belgium is genuinely different from France); Hypothesis~(3) is methodological (Belgium is in principle identifiable like France but with weaker power). The key observation for the joint within-and-across-country pattern in this paper is that the Belgian customs result, taken at face value, says the cross-border leakage channel is small or wrong-signed --- consistent with the within-country result (Sections~\ref{sec:test_h} and~\ref{sec:test_i}) that the across-firm domestic channel is also small. We turn to that joint reading in Section~\ref{sec:discussion}.

\subsection{What we did not run, and why}\label{sec:across:notpursued}

Three Phase~2 and Appendix-B robustness exercises in the original CMdG replication plan were skipped:
\begin{itemize}
    \item \emph{Figure~4 (alternative control with Intrastat-break dummy).} Spec~B in our robustness exercise is the conceptual content of CMdG's Figure~4 (treatment $=$ non-ETS country within the regulated subsample); Spec~B's VCOV failed to converge in the Belgian data, so even if we ran Figure~4 with a 2011 Intrastat-break dummy it would not be identified. Belgium also has no documented analog of the 2011 French Intrastat threshold change, so the break-dummy refinement is unnecessary on substantive grounds.
    \item \emph{Heterogeneity by buyer ETS status (Appendix~B.1).} Splits the buyer firms by their own ETS regulatory status. With our null/negative headline, decomposing it by buyer ETS status only documents which firms are \emph{not} leaking; it does not refine the headline.
    \item \emph{de Chaisemartin--D'Haultfoeuille robustness (Appendix~B.3).} Hedges against negative-weight bias in staggered DiD. Our specification is not staggered (treatment is a fixed cell $\times$ phase interaction), so the negative-weight risk is mechanically zero.
    \item \emph{Leave-one-partner-country-out (Appendix~B.4).} Rules out single big partners (\eg\ China) driving the result. With the headline already null, this only confirms null and adds no information.
\end{itemize}
If the design ever pivots to include Phase~IV (where Belgian carbon prices spiked from $\sim$EUR 30 to EUR 80+ per tCO$_2$), these robustness exercises become useful again because the Phase~IV signal could in principle emerge with a different sign.
